%%%%%%%%%%%%%%%%%%%%%%%%%%%%%%%%%%%%%%%%%%%%%%%%%%%%%
%			HLAVIČKA								%
%%%%%%%%%%%%%%%%%%%%%%%%%%%%%%%%%%%%%%%%%%%%%%%%%%%%%
\documentclass[openany,12pt]{memoir}

\usepackage[utf8]{inputenc} 
\usepackage[czech]{babel}
\usepackage[T1]{fontenc}
\usepackage[top=1.5cm, bottom=2cm, left=2cm, right=2cm]{geometry}  % --> NASTAVENÍ OKRAJŮ
\usepackage{fancyhdr}
\usepackage{graphicx}
\usepackage{lmodern}
\usepackage{xwatermark}
\usepackage{xcolor}
\usepackage{changepage}
\usepackage{newunicodechar}



%%%%%% Package na zpěvník
\usepackage[full]{leadsheets}%http://mirrors.nic.cz/tex-archive/macros/latex/contrib/leadsheets/leadsheets_en.pdf   --> dokumentace	
\definesongtitletemplate{empty}{} 
\setchords{
format = \bfseries,   %tučné akordy
minor = {mi},% 
input-notation = {german},%
output-notation = {german}%
}
\definesongtitletemplate{empty}{} 

\newlength{\drop}
\newwatermark[allpages,color=red!50,angle=0,scale=2, xpos=0,ypos=0]{\includegraphics[width=5cm]{obr/pozadi2.jpg}} %--> dvojka na pozadí


%%%% Vlastní příkazy
\newcounter{Slokočet}   %Automatické číslování slok
\newcommand{\mezera}{\vspace*{0.5cm}}   %Horizontální odsazení slok
\newcommand{\stred}{5.2cm}   %%% Na zarovnání slok doprostřed, pozn. automatičtější zarovnávání na střed nejde
\newcommand{\refren}{\mezera \noindent \textbf{R:} } %refrén
\newcommand{\sloka}{\mezera \noindent \addtocounter{Slokočet}{1} \arabic{Slokočet}. } 	%sloka, která se automaticky čísluje
\newcommand{\ssloka}{\mezera \noindent}  % vlastní číslo sloky
\newcommand{\carka}{,\:}
\newcommand{\m}[1]{\color{white}{#1}}  %Pro akordy

%%%%%%%%%%%%%%%%%%%%%%%%%%%%%%%%%%%%%%%%%%%%%%%%%%%%%%
%			NÁVOD									 %
%%%%%%%%%%%%%%%%%%%%%%%%%%%%%%%%%%%%%%%%%%%%%%%%%%%%%%
%1. Věci v hlavičce IGNOROVAT
%2. Píseň psát do prostoru mezi \begin{song} a \end{song}
%3. další řádek se značí dvěma odsazeníma (= dvakrát stisknout enter)
%4. \refren vždy na začátku refrenu a \sloka na začátek sloky (automaticky se čísluje)
%  \ = alt gr + q ; [/] = alt gr f/g ; {/} = alt gr + b/n; ^ = alt gr + 3 + mezera
%Cokoliv napíšete do ^{  } se bude brát jako akord
%když se toto bude dotýkat nějakého slova (nebude mezi tím a slovem mezera)
%tak se akordy zjeví nad slovem, ale pište to před slovo
%Když se to nedotýká slova, tak akord lítá ve vzduchu a vytiskne se větší mezera
%První možnost je asi preferovanější
%5. Akordy stačí psát jen do první sloky, když se nezmění -- kytaristi to zvládnou



%<++>
%\usepackage{subfiles}
%</++>

\begin{document}

\pagestyle{empty}
%\addtocontents{toc}{\protect\thispagestyle{empty}}
%\tableofcontents \thispagestyle{empty}\newpage
\newgeometry{top=1.5cm, bottom = 0cm, left = 2cm, right = 2cm}




%Pro kompilaci je důležitý, aby ten input a \end{document} zůstali na posledních
%dvou řádcích!!!!
\begin{song}{title=\predtitle\centering Space Oddity \\\large David Bowie  \vspace*{-0.3cm}}  %% sem se napíše jméno songu a autor

\nejvetsi
\begin{centerjustified}

\sloka 
	^{C\z}Ground control to Major ^{Emi\z}Tom, ^{C\z}Ground control to Major ^{Emi\z}Tom:

	^{Ami}Take your ^{Am7/G\z}protein~pills and ^{D7}put your helmet on.

	^{C\z}Ground control to Major ^{Emi\,}Tom:  ^{C\z}Commencing countdown engine's ^{Emi\z}on~~

	^{Ami\z}Check ^{Ami7/G}ignition and may ^{D7\,\,}God's love be with you

	^{C\z }This is ground control to Major ^{E7}Tom, you've really made the ^{F\z}grade!

	And the ^{Fmi\z }papers want to ^{C\z }know whose shirts you ^{F\z}wear,

	now it's ^{Fmi\,}time to leave the ^{C{\color{white}\_\_}}capsule if you ^{F{\color{white}\_}}dare.

	^{C{\color{white}\_}}This is Major Tom to ground ^{\,E7}control, I'm stepping through the ^{F{\color{white}\_}}door

	And I'm ^{Fmi{\color{white}\_}}floating in the ^{C\z}most peculiar ^{F\,\,\,}way 

	and the ^{Fmi\,}stars look very ^{C\z}different ^{F}today

	For ^{Fmaj7}here am I ^{Em7}sitting in a tin ^{Fmaj7\:}can,~far~above the ^{Em7}world

	^{Bmaj7}Planet Earth is ^{Ami}blue and there's ^{G{\color{white}\_\_\_}}nothing I can ^{F}do.

	\phantom{.}

	\phantom{.}

/: \textbf{C\,F\,G\,A} :/ \textbf{Fmaj7\,Em7\,A\,C\,D\,E}

\end{centerjustified}
\newpage
\begin{centerjustified}

\sloka
	^{C{\color{white}\_\_\_}}Though I'm passed one hundred thousand ^{E7{\color{white}\_}}miles, I'm feeling very ^{F\,\,}still

	And I ^{Fmi\,}think my spaceship ^{C\,{\color{white}\_\_}}knows which way to ^{F}go,

	tell my ^{Fmi}wife I love her ^{C\,\,\,}very much, she ^{F{\color{white}\_\_}}knows.

	^{G\z}Ground control to ^{E7\,\,}Major Tom:

	Your ^{Ami{\color{white}\_}}circuit's dead, there's ^{Am7/G\,\,\,}something wrong.

	Can you ^{D7\,\,}hear me Major Tom? Can you ^{C\,\,}hear me Major Tom? 

	Can you ^{G\,\,}hear me Major Tom? Can you\dots

	^{Fmaj7\z}Here~am I ^{Emi7\,\,}floating in my tin ^{Fmaj7\,}can,~far~above the ^{Emi7}moon.

	^{Bmaj7}Planet Earth is ^{Ami}blue and there's ^{G\z}nothing I can ^{F}do.


\end{centerjustified}
\setcounter{Slokočet}{0}
\end{song}

\end{document}
